\originalchapter{覆盖索引、考核工具与数值复现}
\originalsection{冻结范围与正文定位}
本轮在下载前冻结课程与资源清单。主课8.044共22讲幻灯片、四章概率讲义及七份主题笔记；11份作业与答案、两练习卷与答案、两期中与答案、一期末与答案均纳入。进阶仅选8.333讲义15--22与第5组作业；其余讲义、其他作业、历史考试与视频未被本册冒称重编。公开材料共74个PDF、918原始物理页；原始下载SHA256与公开归档SHA256分列记录，五份文件仅分离指定外部图像对象，文字与页数不变。

\begin{longtable}{p{.22\textwidth}p{.68\textwidth}}
\toprule 原材料 & 本册位置及实质范围\\\midrule\endhead
8.044 L1--4、概率Ch1--4 & 第1章：分布、联合条件、物理变量变换、卷积与极限；作业1--3、第一次练习与期中\\
L5--6、系统定义 & 第2章：平衡、温度、功热、响应与全微分；作业3--4\\
L7--9、微正则、间隙原子 & 第3章：Liouville、气体计数、振子边缘、两能级与缺陷；作业4--5、期末1\\
L10--11、Maxwell笔记 & 第4章：四势、稳定性、弹性杆、Carnot；作业6--7、两练习卷\\
L12--13、正则笔记 & 第3、5章：热库、Boltzmann、独立体系与能隙；作业7--8、期中2\\
L14、Raman笔记 & 第6章：正常模、振转、Raman、核自旋；作业8、10\\
L15 & 第7章：辐射模式、Planck、谱峰与状态方程；作业9、期中2\\
L16 & 第8章：Brillouin、磁功、热容、去磁与负温；作业4、6\\
L17、激发态氦 & 第9章：交换、归一化、多体计数与氦；作业10\\
L18 & 第10章：化学势、吸附、巨正则占据；作业10、期末5\\
L19 & 第11章：Fermi海、Sommerfeld、半导体；作业11、期末4\\
L20 & 第12章：引力、白矮星和相对论尺度；作业8\\
L21 & 第13章：声子、零点压力、涨落与氦熔化；作业9、11\\
L22 & 第14章：BEC容量、热容、阱与实验机制；期末5\\
8.333 L15--19 & 第15章：累积量、Mayer、virial、均场、变分和临界；进阶作业1--5\\
L20 & 第6、7、13章：分子、声子与辐射的进阶衔接\\
L21 & 第16章：算符系综、量子熵、自由核与变分；进阶6--10、12\\
L22 & 第9、10、16章：同粒子Hilbert子空间与Fock计数；进阶11\\\bottomrule
\end{longtable}
独立复核不只核算最终值，还逐项核原题条件、小问、原图和正文实质。总数为85道原编号题加作业三一题未编号练习，即86题；270个字母小问和两个嵌套罗马小问。8.044的16份考核对象均有公开官方答案供核对，其中第一次练习答案是扫描页，依题目内容配对；8.333此卷没有公开官方解答，不能用“官方答案”标记。逐题机器清单、独立审阅记录与对象许可清单位于同一源码包的qa目录。

\originalsection{原考试公式表的统一中文版本}
设三个变量由状态方程相联系，局部可逆时有 $x_y|_z=1/(y_x|_z)$、$x_y|_w\,y_z|_w=x_z|_w$、$x_y|_z\,y_z|_x\,z_x|_y=-1$。链式求导不能擅自换固定量；奇点处先用微分约束。

$K$ 个不同对象排列数 $K!$，依序取 $L$ 个为 $K!/(K-L)!$，不计顺序为 $\binom KL$。大数 $\ln K!=K\ln K-K+O(\ln K)$；$\alpha$维球体积 $\pi^{\alpha/2}R^\alpha/\Gamma(\alpha/2+1)$；正边长单纯形体积 $S^d/d!$。

外界对系统的功沿课中约定：流体 $-P\dd V$、线 $\mathcal F\dd L$、膜 $\gamma\dd A$、电介质 $\mathcal E\dd\mathcal P$、磁体 $H\dd M$、可逆电池 $\mathcal E\dd Z$。统一写作 $X\dd x$ 时，$\dd U=T\dd S+X\dd x+\mu\dd N$，$F=U-TS$、$G=U-TS-Xx$、$\mathcal H=U-Xx$，分别微分为 $-S\dd T+X\dd x+\mu\dd N$、$-S\dd T-x\dd X+\mu\dd N$、$T\dd S-x\dd X+\mu\dd N$。巨势 $\Phi=F-\mu N$、$\dd\Phi=-S\dd T-N\dd\mu+X\dd x$。

谐振子 $E_n=(n+1/2)\hbar\omega$，$z=e^{-\beta\hbar\omega/2}/(1-e^{-\beta\hbar\omega})$，$\avg E=\hbar\omega[1/2+(e^{\beta\hbar\omega}-1)^{-1}]$；单粒子平均占据为 $[e^{\beta(\epsilon-\mu)}\pm1]^{-1}$，正号费米、负号玻色。Kirchhoff $e_\omega/\alpha_\omega=cu_\omega/4$，黑体总出射 $e=\sigma_{\rm SB}T^4$；零点项不随温度变，热辐射出射不包含真空零点流。

常用积分（各收敛条件下）为
\begin{align*}
\int_0^\infty x^ne^{-x}\dd x&=n!,&\int_0^\infty x^{-1/2}e^{-x}\dd x&=\sqrt\pi,\\
\int_0^1x^m(1-x)^n\dd x&=\frac{m!n!}{(m+n+1)!},&\int_0^\infty xe^{-x^2}\dd x&=\frac12.
\end{align*}
均值零、标准差 $\sigma$ 的高斯偶矩为 $(2n-1)!!\,\sigma^{2n}$，原部分公式表将指数印作 $\sigma^n$，这里纠正。$\int e^{ax}\dd x=e^{ax}/a$，$\int xe^{ax}\dd x=e^{ax}(ax-1)/a^2$，$\int x^2e^{ax}\dd x=e^{ax}(a^2x^2-2ax+2)/a^3$（$a\ne0$）；$\int\dd x/(1+e^x)=x-\ln(1+e^x)+C$。

常用极限：$\ln n/n\to0$、$n^{1/n}\to1$、$x^{1/n}\to1$（$x>0$）、$x^n\to0$（$|x|<1$）、$(1+x/n)^n\to e^x$、$x^n/n!\to0$（固定 $x$）。这些表格是考核时提供的工具，不构成缺失题面。

\originalsection{实际运行的数值记录}
随包脚本tools/reproduce\_math.py已实跑，输出qa/numerical-reproduction.json。它使用题给常数和已注明的物理常数，以双精度、100次二分、直接收敛谱和复现数值；下表只列关键结果。
\begin{longtable}{p{.44\textwidth}p{.46\textwidth}}
\toprule 复现对象 & 实际结果\\\midrule\endhead
宏观能量相对标准差 & $1.574266415\times10^{-10}$\\
氢第一激发/基态，300/1000 K & $1.776649210\times10^{-171}$；$1.571529363\times10^{-51}$\\
经典核越垒概率 & $2.362720532\times10^{-150}$（使用尾部渐近式）\\
双质子融合事件率 & $1.406988105\times10^{-84}\,\mathrm{s}^{-1}$\\
电容、电感噪声与10MHz功率 & $6.435795988$ 微伏；$2.035177388$ 纳安；$4.141947\times10^{-14}$ W\\
铜/$^3$He 完全极化模型场 & $96316.20350$ T；$5062.628225$ T\\
硬杆全局熵比较 & $\theta_c=1.199177921$；$n_cl^2=0.265565680$\\
硬杆驻值折点 & $\theta_f=1.720667178$；$n_fl^2=0.239623628$\\
转子 $u=.01$ 直接谱和 & $Z=100.334001273$，渐近式 $100.334000000$\\\bottomrule
\end{longtable}
新版另实际枚举300 K转动Raman线的核自旋加权占据数：氮 $I_6=51.972193$、$I_7=26.188730$、$I_8=50.854714$，最强为 $L=6$ 第7条Stokes线；氯 $I_{19}=250.338351$、$I_{20}=150.706093$、$I_{21}=250.832456$，最强为 $L=21$ 第22条。完整参数、各奇偶支峰及枚举范围见JSON的weighted\_rotational\_raman。这里按原notes的占据数模型，核自旋权重不可从单线比较中删去；序数不计Rayleigh中心线和anti-Stokes线。

这些数值验证计算路径与模型中的具体实例。它们不能代替一般证明、临界极限、严格误差界或真实实验数据；转子渐近余项的解析界见进阶作业10，恒星模型和硬杆模型的物理局限在相应题中说明。
